\begin{titlepage}
\thispagestyle{empty}
\begin{tikzpicture}[remember picture,overlay]
\fill[Ink] (current page.south west) rectangle (current page.north east);
\begin{scope}[shift={(current page.south west)},x=1cm,y=1cm]
\fill[white,opacity=.025] (0,10) rectangle (19,18);
\draw[Gold,line width=.7pt] (2,24.2)--(17,24.2);
\draw[Mist,opacity=.32,line width=.5pt] (2,8.4)--(17,8.4);
\foreach \j in {0,...,6}{
 \pgfmathsetmacro{\xx}{2.7+2.15*\j}
 \draw[Mist,opacity=.18,line width=.35pt] (\xx,8.4)--(\xx,16.3);
 \begin{scope}[shift={(\xx,11.5)},xscale=.65,yscale=.7]
  \draw[Gold,opacity=.65,line width=.6pt] (-1.15,-1.25) .. controls (-.3,-.2) and (-.6,.6) .. (0,1.35) .. controls (.6,.6) and (.3,-.2) .. (1.15,-1.25);
  \draw[Mist,opacity=.50,line width=.5pt] (-1.15,.65) .. controls (-.5,.2) and (-.55,-.4) .. (0,-.5) .. controls (.55,-.4) and (.5,.2) .. (1.15,.65);
 \end{scope}
 \fill[Gold] (\xx,8.4) circle (1.15pt);
}
\node[anchor=west,text=Paper,font=\fontsize{11}{14}\selectfont\sffamily] at (2,25.2) {ALEXANDER GROTHENDIECK};
\node[anchor=north west,text=Paper,font=\fontsize{41}{53}\selectfont\bfseries,align=left] at (1.8,22.9) {代数几何\\基础};
\node[anchor=west,text=Gold,font=\fontsize{12}{17}\selectfont\itshape] at (2,17.8) {Fondements de la géométrie algébrique};
\node[anchor=west,text=Mist,font=\small] at (2,16.7) {Bourbaki 报告选集\quad 1957—1962};
\node[anchor=west,text=Paper,font=\fontsize{34}{40}\selectfont\sffamily] at (2,6.9) {FGA};
\node[anchor=west,text=Mist,font=\small] at (2,5.7) {中文译本\quad 附中·英·法术语索引};
\node[anchor=east,text=Mist,font=\small] at (17,2.4) {对偶\enspace ·\enspace 下降\enspace ·\enspace 模空间};
\draw[Gold,line width=.7pt] (2,3.35)--(17,3.35);
\end{scope}
\end{tikzpicture}
\null\vfill
\end{titlepage}
\cleardoublepage
\thispagestyle{empty}
\null\vfill
\begin{flushleft}
{\Large\bfseries\color{Ink}代数几何基础}\par\medskip
Alexander Grothendieck 著\par
\textit{Fondements de la géométrie algébrique}\par
\textit{Extraits du Séminaire Bourbaki 1957--1962}\par\bigskip
原版：Secrétariat mathématique，11 rue Pierre Curie，Paris 5e，1962。\par
底本：用户提供的 208 页扫描本。\par
排版：XeLaTeX；正文与数学公式重新排印，交换图以 TikZ 重绘。\par
本译本为辅助阅读的中文翻译。\par
\end{flushleft}
\vfill
\clearpage
\chapter*{译本说明}
\addcontentsline{toc}{chapter}{译本说明}
本书汇集 Grothendieck 于1957至1962年在 Bourbaki 讨论班所作的八篇报告，以及1962年补写的评论与勘误。译本依底本次序收录全部内容，保留各篇的节号、命题号和公式编号。页边的“原书”标记指向法文印刷页码，例如“190-07”表示第190号报告的第7页；扫描本中149-04有破损页和完整重扫页两张，译文采用完整页，不重复收录。

技术术语在中文后括注法文。点击正文中的术语，可跳至书末中、英、法对照表；表中的各个出现位置可分别跳回正文。为方便和其他文献对读，历史用语 \textit{préschéma} 译作“预概形”，\textit{schéma} 译作“概形”；理解其含义时仍应以本书各处的定义和假设为准。旧称 \textit{simple sur} 的相关用法，按上下文译作“在……上光滑”，与“正则”区别处理。

作者的评论反映1962年时的进展，其中“最近”“目前”“尚未发表”等表述均依原文保留。作者已指出的错误或缺失条件，在相应正文处说明其依据，并链接到书前评论。另加的数学说明以“译者注”标明；这类说明给出补充条件、具体反例或论证，不与原文混排为作者的结论。书末另列译者注索引，便于集中查阅。

本书的报告带有提纲和研究预告的性质；原文未给出的证明，译文不自行写成原作者已完成的证明。文献条目保留原题名，同时补出中文题意，便于检索原文。
